\documentclass[11pt]{article}

\usepackage[margin=1in]{geometry}
\usepackage[T1]{fontenc}
\usepackage[utf8]{inputenc}
\usepackage{lmodern}
\usepackage{amsmath,amssymb}
\usepackage{booktabs,longtable,array}
\usepackage{hyperref}
\usepackage{xcolor}
\usepackage{listings}
\usepackage{enumitem}
\setlength{\emergencystretch}{2em}

\hypersetup{
  colorlinks=true,
  linkcolor=blue!50!black,
  urlcolor=blue!50!black,
  citecolor=blue!50!black
}

\newcolumntype{L}[1]{>{\raggedright\arraybackslash}p{#1}}

\lstdefinestyle{bashstyle}{
  basicstyle=\ttfamily\footnotesize,
  frame=single,
  breaklines=true,
  columns=fullflexible,
  keepspaces=true,
  showstringspaces=false,
  upquote=true
}

\lstdefinestyle{smallbash}{
  basicstyle=\ttfamily\scriptsize,
  frame=single,
  breaklines=true,
  columns=fullflexible,
  keepspaces=true,
  showstringspaces=false,
  upquote=true
}

\title{\texttt{jeener\_packet\_gaussian.py}: Complete User Manual}
\author{Louis-S. Bouchard \and Arianna Wu}
\date{October 10, 2026}

\begin{document}
\maketitle
\noindent\textbf{Numerical status.} The covariance equations do not reproduce
the paper's separable zero-response null control. The code now reports the
computed covariance on that branch. Constraint-tangent injection, multipartite
partitions, physicality checks and repeated-time sampling have been corrected.
The regression values in this manual describe the corrected implementation;
they are not a validation of the published null-control claim. See
\texttt{AUDIT.md} in the repository root for the derivation and limitations.
\tableofcontents

\section{Purpose and Scope}
\texttt{jeener\_packet\_gaussian.py} is an integrated mesoscopic spin-packet
simulator for the eight-spin Jeener cube. It combines deterministic packet
dynamics, optional loop-dressed corrections, stochastic torque, short-memory
damping, native packet-mode Gaussian diagnostics, and an optional
engineered quadratic-Gaussian benchmark layer.

This manual is a standalone reference for the code. It documents the
current executable name, the current command-line interface, the physical
meaning of the main options, the batch commands used for the PRA benchmark
calculations, the expected outputs, and representative regression values.

\section{Runtime Requirements}
The code is a single Python program. It requires a standard scientific Python
environment with:
\begin{itemize}
\item Python 3;
\item NumPy;
\item SciPy;
\item Matplotlib.
\end{itemize}
For noninteractive batch runs, use \texttt{--no-show 1}. If running on a headless
machine, set Matplotlib to an offscreen backend in the shell environment or use a
Python environment already configured for batch plotting.

\section{Program Overview}
\subsection{Physical degrees of freedom}
The native model evolves eight packet variables on the Jeener cube. Each packet
\(j\) has a magnetization-like vector \(\mathbf M_j\) and a response/coherence
vector \(\mathbf N_j\). The initialized packet state is placed on the rotated
coherent-state constraint surface
\begin{equation}
|\mathbf M_j|^2 + |\mathbf N_j|^2 = R^2,
\qquad
\mathbf M_j \cdot \mathbf N_j = 0,
\end{equation}
with \(R=1\) in the default code configuration.

\subsection{Initial-state preparation}
Use \texttt{--Nvec Nx Ny Nz} to specify the response vector. The code projects the
requested response vector into the tangent plane of each packet magnetization
direction and preserves the requested tangent amplitude. The packet
magnetization magnitude is then chosen as
\begin{equation}
|\mathbf M_j| = \sqrt{R^2 - |\mathbf N_j|^2}.
\end{equation}
For the benchmark preparation
\[
\texttt{--Nvec 0.30 0 0},
\]
the initialized packet state has
\[
|\mathbf N_j| = 0.30,
\qquad
|\mathbf M_j| \simeq 0.953939,
\qquad
\mathbf M_j \cdot \mathbf N_j = 0.
\]

\subsection{Main analysis layers}
The program has four main layers:
\begin{enumerate}[label=\arabic*.]
\item spectral and time-domain diagnostics for the deterministic packet model;
\item loop-dressed packet dynamics through \(\alpha\), \(\beta\), \(\sigma\), and
\(\kappa\) coefficients;
\item native packet-mode Gaussian entanglement diagnostics;
\item optional engineered quadratic-Gaussian Hamiltonian benchmarks.
\end{enumerate}

\section{Basic Invocation}
\subsection{General pattern}
The usual batch invocation is:
\begin{lstlisting}[style=bashstyle]
python3 -u ./jeener_packet_gaussian.py [options...]
\end{lstlisting}

For PRA-style native Gaussian heatmaps, the core pattern is:
\begin{lstlisting}[style=bashstyle]
python3 -u ./jeener_packet_gaussian.py \
  --mode eigen \
  --kmax 0.03 \
  --loops 0 \
  --Nvec 0.30 0 0 \
  --tip 0.05 \
  --pattern staggered \
  --ent 1 \
  --ent-save 1 \
  --ent-native-model frozen \
  --ent-native-source full \
  --ent-window 0.0:1.0 \
  --ent-NK 13 \
  --ent-NT 17 \
  --T 8.0 \
  --npts 2401 \
  --fig-dpi 800 \
  --fig-fontsize 12 \
  --fig-cmap magma \
  --fig-transparent 0 \
  --no-show 1
\end{lstlisting}

\subsection{Mode selection}
\begin{longtable}{L{0.22\textwidth}L{0.70\textwidth}}
\toprule
\textbf{Mode} & \textbf{Meaning} \\
\midrule
\endfirsthead
\toprule
\textbf{Mode} & \textbf{Meaning} \\
\midrule
\endhead
\texttt{all} & Runs eigenvalue maps, eigenvalue scatter/stick plots, and one FFT map. This is the default. \\
\texttt{eigen} & Runs eigenvalue color maps, scatter plots, and stick spectra. This is the standard mode for the PRA entanglement command suite. \\
\texttt{map} & Runs FFT color maps over a \(k\)-sweep. \\
\texttt{fft} & Runs a per-spin FFT at \(k=0\). \\
\texttt{growth} & Runs the maximum real eigenvalue growth diagnostic versus \(k\). \\
\bottomrule
\end{longtable}

\section{Native Packet-Gaussian Entanglement}
\subsection{Computed native models}
The native entanglement analysis is enabled with \texttt{--ent 1}. The key switch
is \texttt{--ent-native-model}:
\begin{longtable}{L{0.24\textwidth}L{0.68\textwidth}}
\toprule
\textbf{Model} & \textbf{Meaning} \\
\midrule
\endfirsthead
\toprule
\textbf{Model} & \textbf{Meaning} \\
\midrule
\endhead
\texttt{frozen} & Uses the initial projected native packet generator. This is the benchmark model used for the PRA reproduction commands in this manual. \\
\texttt{linearized} & Propagates the native Gaussian covariance with a time-dependent tangent-map model obtained by finite-differencing the nonlinear one-step map around the deterministic background trajectory. This is slower and useful for validation. \\
\texttt{none} & Disables native covariance output; no entanglement result is reported. \\
\bottomrule
\end{longtable}

\subsection{Invariant branch}
When \(\mathbf N(0)=0\), the implemented native branch preserves
\(\mathbf N(t)=0\). The dipolar magnetization fluctuation generator is still
nonzero, so this response invariant does not imply separable covariance.
The command below computes that covariance and warns that the paper's null
control is not reproduced. No separable baseline is imposed:
\begin{lstlisting}[style=bashstyle]
python3 -u ./jeener_packet_gaussian.py \
  --mode eigen --kmax 0.03 --loops 0 --Nvec 0 0 0 \
  --ent 1 --ent-native-model frozen --ent-native-source full \
  --ent-NK 4 --ent-NT 9 --T 8.0 --npts 801 --no-show 1
\end{lstlisting}

\subsection{Pairwise and multipartite outputs}
The native Gaussian analysis can compute pairwise metrics, multipartite
bipartition metrics, or both. The standard PRA commands use
\texttt{--ent-layer both}, which is the default. Pair selection is controlled by
\texttt{--ent-pairs}; multipartite subset selection is controlled by
\texttt{--ent-multi-size}, \texttt{--ent-multi-subsets}, and
\texttt{--ent-multi-partitions}.

Each three-packet subset is tested across all three one-versus-two
bipartitions. The score for full inseparability uses the weakest partition, then
the requested aggregation over subsets. The default aggregation is
\texttt{max}; \texttt{mean} averages both pair and subset metrics.

The physical packet perturbation includes
\(\delta\mathbf N_j=-\mathbf M_j(\mathbf N_j\cdot\delta\mathbf M_j)/|\mathbf M_j|^2\),
which preserves both linearized packet constraints. A canonical basis check
alone does not certify the Gaussian evolution. The code also checks covariance
uncertainty and Gaussian complete positivity. Frozen evolution advances
directly between requested sample times with exact Gaussian steps.

\section{Loop-Dressed Physics Controls}
The loop order is selected with \texttt{--loops}. The allowed range is 0 through
4. Nonzero loop order only changes the calculation when one or more loop
coefficient lists are supplied.

\begin{longtable}{L{0.20\textwidth}L{0.72\textwidth}}
\toprule
\textbf{Coefficient} & \textbf{Role} \\
\midrule
\endfirsthead
\toprule
\textbf{Coefficient} & \textbf{Role} \\
\midrule
\endhead
\texttt{--alpha} & Local coherent frequency-shift coefficients. These retune local packet precession and do not by themselves provide the dominant nonlocal entangling enhancement in the benchmark scans. \\
\texttt{--beta} & Dissipative damping coefficients. These are typically suppressive for Gaussian entanglement diagnostics. \\
\texttt{--sigma} & Coherent renormalization of the nonlocal dipolar \(M\)-\(M\) kernel. This is the main coherent enhancement channel in the benchmark hierarchy. \\
\texttt{--kappa} & Coherent back-action renormalization. This can contribute with \(\sigma\), but is weaker by itself in the benchmark scans. \\
\bottomrule
\end{longtable}

\section{Output Files}
\subsection{Figure and data routing}
By default, figures and exported \texttt{.npz} sidecar files are written to the
current working directory. Use \texttt{--outdir <directory>} to route outputs to a
dedicated folder. Use \texttt{--export-data 1} to export data sidecars when the
code has a data payload available for the plot.

\subsection{Standard native entanglement heatmaps}
With \texttt{--ent-save 1}, the native Gaussian analysis writes:
\begin{lstlisting}[style=smallbash]
ent_pairwise_gaussian_logneg_heatmap.png
ent_pairwise_gaussian_nupt_heatmap.png
ent_pairwise_gaussian_duan_heatmap.png
ent_multipartite_lnmin_heatmap.png
ent_multipartite_full_insep_heatmap.png
\end{lstlisting}
For large command suites, append run-specific suffixes using a shell helper such
as the \texttt{tag\_ent} function in Section~\ref{sec:pra-suite}.

\subsection{Spectral and FFT outputs}
The spectral and FFT layers save filenames containing the \(N\)-vector tag,
\(k_{\max}\), and loop order. Typical outputs include:
\begin{lstlisting}[style=smallbash]
eigen_map_Nvec_<tag>_kmax_<value>_loops_<order>.png
eigen_scatter_Nvec_<tag>_kmax_<value>_loops_<order>.png
eigen_sticks_k_<value>_Nvec_<tag>_loops_<order>.png
fft_map_Nvec_<tag>_kmax_<value>_loops_<order>_<kind>.png
\end{lstlisting}

\section{Command-Line Option Reference}
\subsection{Execution, time, and spectral controls}
\begin{longtable}{L{0.28\textwidth}L{0.64\textwidth}}
\toprule
\textbf{Option} & \textbf{Use} \\
\midrule
\endfirsthead
\toprule
\textbf{Option} & \textbf{Use} \\
\midrule
\endhead
\texttt{--mode} \{all, eigen, map, fft, growth\} & Selects the diagnostic layer to run. Default: \texttt{all}. \\
\texttt{--k} & Representative \(k\) value for the eigen stick spectrum. Defaults to \(0.8 k_{\max}\). \\
\texttt{--kmax} & Maximum \(k\) used for swept plots. Default in noninteractive mode: \texttt{0.1}. \\
\texttt{--NK} & Number of \(k\) samples for FFT maps; also used as a convenient override in some swept diagnostics. \\
\texttt{--T} & Time-domain duration. Default: \texttt{200.0}. \\
\texttt{--npts} & Number of time samples. Default: \texttt{8001}. \\
\texttt{--tip} & Initial packet magnetization tip angle in radians. Default: \texttt{0.05}. \\
\texttt{--grange} & Optional symmetric vertical limit for the growth plot. \\
\texttt{--show-growth} & When \texttt{--mode all}, also generate the growth curve if set to 1. \\
\texttt{--fft-norm} & Normalize each \(k\)-row in FFT maps. Default: \texttt{1}. \\
\texttt{--frange} & Frequency half-window for FFT and spectral plots. Default: \texttt{0.25}. \\
\texttt{--plotsigma} & Gaussian smear width for FFT maps. Default: \texttt{0.0015}. \\
\texttt{--NFFT} & FFT zero-padding length. Default: \texttt{131072}. \\
\texttt{--interactive} & If set to 1, prompts for omitted parameters in a terminal. Default: \texttt{0}. \\
\bottomrule
\end{longtable}

\subsection{Initial state controls}
\begin{longtable}{L{0.28\textwidth}L{0.64\textwidth}}
\toprule
\textbf{Option} & \textbf{Use} \\
\midrule
\endfirsthead
\toprule
\textbf{Option} & \textbf{Use} \\
\midrule
\endhead
\texttt{--Nvec Nx Ny [Nz]} & Preferred interface for the base response vector. If \(N_z\) is omitted it is taken as 0. \\
\texttt{--N0} & Compatibility interface for setting response magnitude. Prefer \texttt{--Nvec}. \\
\texttt{--n\_hat} & Direction used with \texttt{--N0}. \\
\texttt{--pattern} \{staggered, uniform, noisy\} & Distribution pattern for the base response over the cube. The PRA benchmarks use \texttt{staggered}. \\
\bottomrule
\end{longtable}

\subsection{Loop and nonlinear controls}
\begin{longtable}{L{0.28\textwidth}L{0.64\textwidth}}
\toprule
\textbf{Option} & \textbf{Use} \\
\midrule
\endfirsthead
\toprule
\textbf{Option} & \textbf{Use} \\
\midrule
\endhead
\texttt{--loops} & Loop correction order, clamped to 0 through 4. Default in noninteractive mode: \texttt{0}. \\
\texttt{--alpha} & \(\delta\omega/k\) coefficient list \(\alpha_1,\ldots,\alpha_4\). \\
\texttt{--beta} & \(\lambda_\perp/k\) coefficient list \(\beta_1,\ldots,\beta_4\). \\
\texttt{--sigma} & Coherent dipolar-kernel coefficient list \(\sigma_1,\ldots,\sigma_4\). \\
\texttt{--kappa} & Coherent back-action coefficient list \(\kappa_2,\ldots,\kappa_4\). \\
\texttt{--alpha-k2} & Extra local-shift scaling by \(1+a_{k2}k^2\). Default: \texttt{0.0}. \\
\texttt{--beta-k2} & Extra damping scaling by \(1+b_{k2}k^2\). Default: \texttt{0.0}. \\
\texttt{--beta-anis} & Damping anisotropy, with \(\lambda_x=\lambda(1+\mathrm{ba})\) and \(\lambda_y=\lambda(1-\mathrm{ba})\). Default: \texttt{0.0}. \\
\texttt{--dynamic-mix} & Use \(k|\mathbf N_a(t)|\) in the \(N\)-mixing field when 1; use \(k|\mathbf N_0|\) when 0. Default: \texttt{1}. \\
\texttt{--cubic-eta} & Explicit \(|N|^2 dN\) inertial correction. Default: \texttt{0.0}. \\
\texttt{--prefer-cubic} & If 1, use explicit cubic correction and bypass the multiplicative vertex term. Default: \texttt{0}. \\
\bottomrule
\end{longtable}

\subsection{Noise, memory, and metrics}
\begin{longtable}{L{0.28\textwidth}L{0.64\textwidth}}
\toprule
\textbf{Option} & \textbf{Use} \\
\midrule
\endfirsthead
\toprule
\textbf{Option} & \textbf{Use} \\
\midrule
\endhead
\texttt{--noise} & Enable stochastic torque. Default in noninteractive mode: \texttt{0}. \\
\texttt{--sde-scheme \{ito,strat\}} & Selects Ito/Euler-Maruyama or Stratonovich/Heun integration. Default: \texttt{ito}. \\
\texttt{--noise-scale} & FDT-like scale \(D_\perp=\mathrm{noise\_scale}\lambda_\perp\). Default: \texttt{0.0}. \\
\texttt{--noise-target \{M,MN\}} & Apply stochastic torque to \(M\) only or to both \(M\) and \(N\). Default: \texttt{M}. \\
\texttt{--nreal} & Number of stochastic realizations for FFT averaging. Default: \texttt{1}. \\
\texttt{--seed} & Random-number seed. Default: \texttt{12345}. \\
\texttt{--mem} & Enable short-memory damping through a Prony expansion. Default in noninteractive mode: \texttt{0}. \\
\texttt{--mem-target \{M,MN\}} & Apply memory to \(M\) only or to both \(M\) and \(N\). Default: \texttt{M}. \\
\texttt{--mem-L} & Number of memory exponentials. Default: \texttt{0}. \\
\texttt{--mem-tau}, \texttt{--mem-a} & Time constants and amplitudes for the Prony memory kernel. \\
\texttt{--mem-tangent} & Apply memory only to transverse components when 1. Default: \texttt{1}. \\
\texttt{--metrics} & Print deterministic summary metrics after figure generation. Default: \texttt{0}. \\
\texttt{--metrics-k} & Fraction of \(k_{\max}\) used for the metrics run. Default: \texttt{0.8}. \\
\bottomrule
\end{longtable}

\subsection{Native Gaussian entanglement options}
\begin{longtable}{L{0.31\textwidth}L{0.61\textwidth}}
\toprule
\textbf{Option} & \textbf{Use} \\
\midrule
\endfirsthead
\toprule
\textbf{Option} & \textbf{Use} \\
\midrule
\endhead
\texttt{--ent} & Enable entanglement analysis. Default: \texttt{0}. \\
\texttt{--ent-layer} \{pair, multi, both, gaussian\} & Select pairwise, multipartite, both, or the pairwise Gaussian alias. Default: \texttt{both}. \\
\texttt{--ent-pairs \{nn,all,custom\}} & Pair set for pairwise Gaussian metrics. Default: \texttt{nn}. \\
\texttt{--ent-custom-pairs} & Custom pair list such as \texttt{0-1,2-5}. Used with \texttt{--ent-pairs custom}. \\
\texttt{--ent-window} & Time-window fraction \(t_0:t_1\) sampled for heatmaps. Default: \texttt{0.8:1.0}. \\
\texttt{--ent-S} & Effective packet spin length used for Holstein-Primakoff normalization. Default: \texttt{100.0}. \\
\texttt{--ent-aggregate \{max,mean\}} & Aggregate selected pairs/subsets by best case or mean. Default: \texttt{max}. \\
\texttt{--ent-NK} & Number of \(k\) samples for entanglement heatmaps. Defaults to \texttt{--NK} or 31. \\
\texttt{--ent-NT} & Number of time bins inside the chosen window. Default: \texttt{41}. \\
\texttt{--ent-save} & Save entanglement heatmap figures. Default: \texttt{1}. \\
\texttt{--ent-cv-rot} & If positive, optimize the pairwise Duan witness over a local phase-rotation grid. Default: \texttt{0}. \\
\texttt{--ent-phys-regularize} & Add the smallest isotropic noise needed to make sampled covariances quantum-physical before evaluating Gaussian entanglement. Default: \texttt{1}. \\
\texttt{--ent-multi-size} & Subset size for multipartite analysis when no custom subsets are supplied. Default: \texttt{3}. \\
\texttt{--ent-multi-subsets} & Custom multipartite subsets separated by semicolons, e.g. \texttt{0-1-2;0-1-4-5}. \\
\texttt{--ent-multi-partitions \{1vr,all\}} & Test one-vs-rest or all unique bipartitions inside each selected subset. Default: \texttt{all}. \\
\texttt{--ent-native-tol} & Tolerance for detecting the exact native invariant branch \(N(t)=0\). Default: \texttt{1e-15}. \\
\texttt{--ent-native-model \{linearized,frozen,none\}} & Select the native model. Default: \texttt{linearized}. Benchmark commands use \texttt{frozen}. \\
\texttt{--ent-native-source \{full,diag,cross,mm-only\}} & Diagnostic block-selection switch. In the canonical packet-mode model, the physical full tangent map is used. Default: \texttt{full}. \\
\texttt{--ent-native-nth} & Initial thermal occupancy per native Gaussian mode. Default: \texttt{0.0}. \\
\texttt{--ent-native-nbath} & Bath thermal occupancy for explicit native local damping terms. Default: \texttt{0.0}. \\
\texttt{--ent-native-reg-warn} & Warning threshold for non-negligible native covariance repair. Default: \texttt{1e-6}. \\
\bottomrule
\end{longtable}

\subsection{Plotting and export options}
\begin{longtable}{L{0.28\textwidth}L{0.64\textwidth}}
\toprule
\textbf{Option} & \textbf{Use} \\
\midrule
\endfirsthead
\toprule
\textbf{Option} & \textbf{Use} \\
\midrule
\endhead
\texttt{--fig-dpi} & DPI for saved figures. Default: \texttt{300}. \\
\texttt{--fig-cmap} & Default heatmap colormap. Default: \texttt{magma}. \\
\texttt{--fig-fontsize} & Base Matplotlib font size. Default: \texttt{11.0}. \\
\texttt{--fig-transparent} & Save transparent figure backgrounds when 1. Default: \texttt{0}. \\
\texttt{--no-show} & Suppress interactive \texttt{plt.show()} calls. Use \texttt{1} for batch runs. \\
\texttt{--outdir} & Optional output directory for figures and exported data. \\
\texttt{--export-data} & Export \texttt{.npz} sidecar data when available. Default: \texttt{1}. \\
\bottomrule
\end{longtable}

\subsection{Engineered Gaussian Hamiltonian options}
\begin{longtable}{L{0.31\textwidth}L{0.61\textwidth}}
\toprule
\textbf{Option} & \textbf{Use} \\
\midrule
\endfirsthead
\toprule
\textbf{Option} & \textbf{Use} \\
\midrule
\endhead
\texttt{--eng} & Enable engineered Gaussian Hamiltonian analysis. Default: \texttt{0}. \\
\texttt{--eng-preset} & Convenience preset: \texttt{none}, \texttt{cube-tms}, \texttt{cube-cluster}, \texttt{graph-cz}, \texttt{tact}, \texttt{tact-uniform}, \texttt{tact-checker}, \texttt{oatq}, \texttt{oatp}, or \texttt{hybrid}. \\
\texttt{--eng-graph} & Graph for engineered edge terms: \texttt{cube-nn}, \texttt{cube}, \texttt{line}, \texttt{ring}, \texttt{star}, \texttt{complete}, \texttt{custom}, or \texttt{none}. \\
\texttt{--eng-edges} & Custom weighted edge list such as \texttt{0-1:1,1-2:0.5}. \\
\texttt{--eng-graph-weight} & Overall graph-edge scale. Default: \texttt{1.0}. \\
\texttt{--eng-bs} & Beam-splitter edge coefficient. Default: \texttt{0.0}. \\
\texttt{--eng-tms} & Two-mode-squeezing edge coefficient. Default: \texttt{0.0}. \\
\texttt{--eng-cz} & Controlled-\(Z\)/\(Q_iQ_j\) graph interaction coefficient. Default: \texttt{0.0}. \\
\texttt{--eng-sms} & Local single-mode squeezing coefficient. Default: \texttt{0.0}. \\
\texttt{--eng-tact} & Collective TACT-like squeezing coefficient. Default: \texttt{0.0}. \\
\texttt{--eng-oatq}, \texttt{--eng-oatp} & Collective OAT-like \(Q_u^2\) or \(P_u^2\) coefficients. Defaults: \texttt{0.0}. \\
\texttt{--eng-omega} & Local harmonic rotation frequency. Default: \texttt{0.0}. \\
\texttt{--eng-mode} & Collective supermode for engineered OAT/TACT terms. Choices include \texttt{uniform}, \texttt{bright}, \texttt{checker}, \texttt{staggered}, axis/layer modes, and \texttt{custom}. \\
\texttt{--eng-mode-custom} & Eight custom weights for \texttt{--eng-mode custom}. \\
\texttt{--eng-init} & Initial engineered Gaussian state: \texttt{coherent}, \texttt{vacuum}, \texttt{thermal}, \texttt{p-sq}, or \texttt{q-sq}. \\
\texttt{--eng-r} & Single-mode squeezing parameter for \texttt{p-sq} or \texttt{q-sq}. Default: \texttt{0.0}. \\
\texttt{--eng-nth-init} & Initial thermal occupancy for \texttt{--eng-init thermal}. Default: \texttt{0.0}. \\
\texttt{--eng-gamma}, \texttt{--eng-nbath} & Uniform linear damping rate and bath occupancy. Defaults: \texttt{0.0}. \\
\texttt{--eng-profile} & Interaction envelope: \texttt{constant}, \texttt{ramp}, \texttt{pulse}, \texttt{ramphold}, \texttt{rampholdquench}, \texttt{sin2}, or \texttt{echo}. \\
\texttt{--eng-ramp-frac}, \texttt{--eng-hold-frac} & Ramp and hold fractions for ramp-hold profiles. Defaults: \texttt{0.2}, \texttt{0.8}. \\
\texttt{--eng-pulse-on}, \texttt{--eng-pulse-off} & Pulse turn-on and turn-off fractions. Defaults: \texttt{0.0}, \texttt{1.0}. \\
\texttt{--eng-echo-frac} & Echo sign-flip fraction. Default: \texttt{0.5}. \\
\texttt{--eng-T}, \texttt{--eng-npts} & Engineered evolution time and time samples. Defaults: \texttt{40.0}, \texttt{801}. \\
\texttt{--eng-window} & Time-window fraction for engineered metric heatmaps. Default: \texttt{0.0:1.0}. \\
\texttt{--eng-gmax}, \texttt{--eng-NG}, \texttt{--eng-NT} & Engineered interaction sweep maximum, number of interaction samples, and number of time bins. Defaults: \texttt{1.0}, \texttt{21}, \texttt{41}. \\
\texttt{--eng-full} & Compute full eight-mode bipartition metrics. Default: \texttt{1}. \\
\texttt{--eng-save} & Save engineered-analysis figures. Default: \texttt{1}. \\
\texttt{--eng-reg-warn} & Warning threshold for engineered covariance regularization. Default: \texttt{1e-6}. \\
\bottomrule
\end{longtable}

\section{PRA Benchmark Command Suite}
\label{sec:pra-suite}
The commands below reproduce the standard native packet-Gaussian benchmark
suite using the packaged executable.

\subsection{Batch helper}
\begin{lstlisting}[style=bashstyle]
tag_ent () {
  s="$1"
  mv -f ent_pairwise_gaussian_logneg_heatmap.png  "ent_pairwise_gaussian_logneg_heatmap_${s}.png"
  mv -f ent_pairwise_gaussian_nupt_heatmap.png    "ent_pairwise_gaussian_nupt_heatmap_${s}.png"
  mv -f ent_pairwise_gaussian_duan_heatmap.png    "ent_pairwise_gaussian_duan_heatmap_${s}.png"
  mv -f ent_multipartite_lnmin_heatmap.png        "ent_multipartite_lnmin_heatmap_${s}.png"
  mv -f ent_multipartite_full_insep_heatmap.png   "ent_multipartite_full_insep_heatmap_${s}.png"
}

CODE=./jeener_packet_gaussian.py
\end{lstlisting}

\subsection{Classical invariant-manifold control}
\begin{lstlisting}[style=bashstyle]
python3 -u "$CODE" \
  --mode eigen \
  --kmax 0.03 \
  --loops 0 \
  --Nvec 0 0 0 \
  --ent 1 \
  --ent-save 1 \
  --ent-native-model frozen \
  --ent-native-source full \
  --ent-NK 4 \
  --ent-NT 9 \
  --T 8.0 \
  --npts 801 \
  --fig-dpi 400 \
  --fig-fontsize 11 \
  --fig-cmap magma \
  --fig-transparent 0 \
  --no-show 1 | tee classical_control_x000_T8_jeener_packet_gaussian.log
tag_ent classical_control_x000_T8_jeener_packet_gaussian
\end{lstlisting}

\subsection{Native benchmark}
\begin{lstlisting}[style=bashstyle]
python3 -u "$CODE" \
  --mode eigen \
  --kmax 0.03 \
  --loops 0 \
  --Nvec 0.30 0 0 \
  --tip 0.05 \
  --pattern staggered \
  --ent 1 \
  --ent-save 1 \
  --ent-native-model frozen \
  --ent-native-source full \
  --ent-window 0.0:1.0 \
  --ent-NK 13 \
  --ent-NT 17 \
  --T 8.0 \
  --npts 2401 \
  --fig-dpi 800 \
  --fig-fontsize 12 \
  --fig-cmap magma \
  --fig-transparent 0 \
  --no-show 1 | tee fig3_native_corr_x030_T8_jeener_packet_gaussian.log
tag_ent fig3_native_corr_x030_T8_jeener_packet_gaussian
\end{lstlisting}

\subsection{Coherent \texorpdfstring{\(\sigma\)}{sigma} enhancement}
\begin{lstlisting}[style=bashstyle]
python3 -u "$CODE" \
  --mode eigen \
  --kmax 0.03 \
  --loops 4 \
  --sigma 2.00 5.00 10.00 \
  --Nvec 0.30 0 0 \
  --tip 0.05 \
  --pattern staggered \
  --ent 1 \
  --ent-save 1 \
  --ent-native-model frozen \
  --ent-native-source full \
  --ent-window 0.0:1.0 \
  --ent-NK 13 \
  --ent-NT 17 \
  --T 8.0 \
  --npts 2401 \
  --fig-dpi 800 \
  --fig-fontsize 12 \
  --fig-cmap magma \
  --fig-transparent 0 \
  --no-show 1 | tee fig4_sigma_x030_T8_jeener_packet_gaussian.log
tag_ent fig4_sigma_x030_T8_jeener_packet_gaussian
\end{lstlisting}

\subsection{Stronger coherent regime}
\begin{lstlisting}[style=bashstyle]
python3 -u "$CODE" \
  --mode eigen \
  --kmax 0.03 \
  --loops 4 \
  --alpha 5.00 5.00 5.00 5.00 \
  --sigma 5.00 10.00 20.00 \
  --kappa 5.00 10.00 20.00 \
  --Nvec 0.30 0 0 \
  --tip 0.05 \
  --pattern staggered \
  --ent 1 \
  --ent-save 1 \
  --ent-native-model frozen \
  --ent-native-source full \
  --ent-window 0.0:1.0 \
  --ent-NK 13 \
  --ent-NT 17 \
  --T 8.0 \
  --npts 2401 \
  --fig-dpi 800 \
  --fig-fontsize 12 \
  --fig-cmap magma \
  --fig-transparent 0 \
  --no-show 1 | tee fig5_aggr_x030_T8_jeener_packet_gaussian.log
tag_ent fig5_aggr_x030_T8_jeener_packet_gaussian
\end{lstlisting}

\subsection{Thermal suppression}
\begin{lstlisting}[style=bashstyle]
python3 -u "$CODE" \
  --mode eigen \
  --kmax 0.03 \
  --loops 0 \
  --Nvec 0.30 0 0 \
  --tip 0.05 \
  --pattern staggered \
  --ent 1 \
  --ent-save 1 \
  --ent-native-model frozen \
  --ent-native-source full \
  --ent-native-nth 0.10 \
  --ent-window 0.0:1.0 \
  --ent-NK 13 \
  --ent-NT 17 \
  --T 8.0 \
  --npts 2401 \
  --fig-dpi 800 \
  --fig-fontsize 12 \
  --fig-cmap magma \
  --fig-transparent 0 \
  --no-show 1 | tee app_thermal_nth010_x030_T8_jeener_packet_gaussian.log
tag_ent app_thermal_nth010_x030_T8_jeener_packet_gaussian
\end{lstlisting}

\subsection{Appendix control scans}
\begin{lstlisting}[style=bashstyle]
# sigma + kappa
python3 -u "$CODE" --mode eigen --kmax 0.03 --loops 4 \
  --sigma 2.00 5.00 10.00 --kappa 2.00 5.00 10.00 \
  --Nvec 0.30 0 0 --tip 0.05 --pattern staggered \
  --ent 1 --ent-save 1 --ent-native-model frozen --ent-native-source full \
  --ent-window 0.0:1.0 --ent-NK 4 --ent-NT 9 \
  --T 8.0 --npts 1601 --fig-dpi 600 --fig-fontsize 11 \
  --fig-cmap magma --fig-transparent 0 --no-show 1 \
  | tee app_sigma_kappa_x030_T8_jeener_packet_gaussian.log
tag_ent app_sigma_kappa_x030_T8_jeener_packet_gaussian

# alpha only
python3 -u "$CODE" --mode eigen --kmax 0.03 --loops 4 \
  --alpha 5.00 5.00 5.00 5.00 \
  --Nvec 0.30 0 0 --tip 0.05 --pattern staggered \
  --ent 1 --ent-save 1 --ent-native-model frozen --ent-native-source full \
  --ent-window 0.0:1.0 --ent-NK 4 --ent-NT 9 \
  --T 8.0 --npts 1601 --fig-dpi 600 --fig-fontsize 11 \
  --fig-cmap magma --fig-transparent 0 --no-show 1 \
  | tee app_alpha_only_x030_T8_jeener_packet_gaussian.log
tag_ent app_alpha_only_x030_T8_jeener_packet_gaussian

# kappa only
python3 -u "$CODE" --mode eigen --kmax 0.03 --loops 4 \
  --kappa 2.00 5.00 10.00 \
  --Nvec 0.30 0 0 --tip 0.05 --pattern staggered \
  --ent 1 --ent-save 1 --ent-native-model frozen --ent-native-source full \
  --ent-window 0.0:1.0 --ent-NK 4 --ent-NT 9 \
  --T 8.0 --npts 1601 --fig-dpi 600 --fig-fontsize 11 \
  --fig-cmap magma --fig-transparent 0 --no-show 1 \
  | tee app_kappa_only_x030_T8_jeener_packet_gaussian.log
tag_ent app_kappa_only_x030_T8_jeener_packet_gaussian

# beta only
python3 -u "$CODE" --mode eigen --kmax 0.03 --loops 4 \
  --beta 1.00 2.00 4.00 8.00 --beta-anis 0.50 \
  --Nvec 0.30 0 0 --tip 0.05 --pattern staggered \
  --ent 1 --ent-save 1 --ent-native-model frozen --ent-native-source full \
  --ent-window 0.0:1.0 --ent-NK 4 --ent-NT 9 \
  --T 8.0 --npts 1601 --fig-dpi 600 --fig-fontsize 11 \
  --fig-cmap magma --fig-transparent 0 --no-show 1 \
  | tee app_beta_only_x030_T8_jeener_packet_gaussian.log
tag_ent app_beta_only_x030_T8_jeener_packet_gaussian
\end{lstlisting}

\subsection{Optional time-dependent validation runs}
\begin{lstlisting}[style=bashstyle]
python3 -u "$CODE" \
  --mode eigen --kmax 0.03 --loops 0 \
  --Nvec 0.30 0 0 --tip 0.05 --pattern staggered \
  --ent 1 --ent-save 1 \
  --ent-native-model linearized --ent-native-source full \
  --ent-window 0.0:1.0 --ent-NK 9 --ent-NT 17 \
  --T 8.0 --npts 2401 --fig-dpi 800 --fig-fontsize 12 \
  --fig-cmap magma --fig-transparent 0 --no-show 1 \
  | tee app_td_native_x030_T8_jeener_packet_gaussian.log
tag_ent app_td_native_x030_T8_jeener_packet_gaussian

python3 -u "$CODE" \
  --mode eigen --kmax 0.03 --loops 4 \
  --sigma 2.00 5.00 10.00 \
  --Nvec 0.30 0 0 --tip 0.05 --pattern staggered \
  --ent 1 --ent-save 1 \
  --ent-native-model linearized --ent-native-source full \
  --ent-window 0.0:1.0 --ent-NK 9 --ent-NT 17 \
  --T 8.0 --npts 2401 --fig-dpi 800 --fig-fontsize 12 \
  --fig-cmap magma --fig-transparent 0 --no-show 1 \
  | tee app_td_sigma_x030_T8_jeener_packet_gaussian.log
tag_ent app_td_sigma_x030_T8_jeener_packet_gaussian
\end{lstlisting}

\section{Representative Regression Values}
Table~\ref{tab:regression} records representative endpoint values for the
standard frozen native packet-Gaussian benchmark suite. These are useful for
checking that a local run is in the expected numerical regime.

\begin{table}[ht]
\centering
\setlength{\tabcolsep}{3pt}
\begin{tabular}{L{0.33\textwidth}L{0.20\textwidth}L{0.20\textwidth}L{0.19\textwidth}}
\toprule
\textbf{Scenario} & \textbf{Pairwise endpoint} & \textbf{Multipartite endpoint} & \textbf{Comment} \\
\midrule
Invariant branch, \(\mathbf N=0\), loops \(=0\) & \(E_N \simeq 0.7111\) & \(FI \simeq 0.1908\) & Paper null not reproduced. \\
Native benchmark, \(\mathbf N=(0.30,0,0)\), loops \(=0\) & \(E_N \simeq 0.6797\) & \(FI \simeq 0.1873\) & Computed result. \\
\(\sigma\)-enhanced coherent run & \(E_N \simeq 0.7696\) & \(FI \simeq 0.2096\) & Computed result. \\
Stronger coherent regime & \(E_N \simeq 0.8400\) & \(FI \simeq 0.2333\) & Computed result. \\
Thermal run, \texttt{--ent-native-nth 0.10} & \(E_N \simeq 0.4167\) & \(FI \simeq 0.1248\) & Computed result. \\
\bottomrule
\end{tabular}
\caption{Representative high-\(k\), end-of-window values for the benchmark
heatmaps at \(k=0.03\), \(t=8\), tip \(0.05\), staggered pattern and maximum
aggregation. These are numerical regression values, not paper validation.}
\label{tab:regression}
\end{table}

\section{Troubleshooting and Practical Notes}
\begin{itemize}
\item If \texttt{--loops} is nonzero but no loop coefficients are supplied, the code
prints a warning because the loop-dressed spectra and entanglement diagnostics
will look unchanged.
\item For scripted runs, always use \texttt{--no-show 1}. This prevents Matplotlib
from blocking execution.
\item For reproducible output names, either run each command in a clean output
directory with \texttt{--outdir}, or use the \texttt{tag\_ent} helper immediately
after each command.
\item Use \texttt{--Nvec} for new calculations. The magnitude/direction interface
\texttt{--N0}/\texttt{--n\_hat} remains available for compatibility, but is less
explicit.
\item The \texttt{linearized} native model is slower than \texttt{frozen}; use it for
validation or sensitivity checks rather than for the fastest benchmark sweep.
\item Covariance repair does not validate the underlying quantum channel.
Complete-positivity failures exclude affected bins. Inspect all validity
diagnostics and the accepted sample count before interpreting a summary.
\end{itemize}

\section{Minimal Smoke Tests}
The following checks are quick ways to confirm that the local executable is
working:
\begin{lstlisting}[style=bashstyle]
python3 ./jeener_packet_gaussian.py --help

python3 -u ./jeener_packet_gaussian.py \
  --mode eigen --kmax 0.03 --loops 0 --Nvec 0 0 0 \
  --ent 1 --ent-native-model frozen --ent-NK 2 --ent-NT 3 \
  --T 1.0 --npts 101 --no-show 1

python3 -u ./jeener_packet_gaussian.py \
  --mode eigen --kmax 0.03 --loops 0 --Nvec 0.30 0 0 \
  --ent 1 --ent-native-model frozen --ent-NK 2 --ent-NT 3 \
  --T 1.0 --npts 101 --no-show 1
\end{lstlisting}

\end{document}
